\oldpage[IV-2]{36}
\section{Ciclos primos asociados y descomposiciones primarias}

En este parágrafo damos sobre todo la traducción de los resultados sobre los
módulos expuestos en Bourbaki, \emph{Alg. comm.}, cap.~IV, que seguimos muy de
cerca. Las nociones que siguen parecen tener poco interés salvo en el caso de
preesquemas \emph{localmente noetherianos}.

\subsection{Ciclos primos asociados a un módulo.}

\begin{definition}[3.1.1]
\label{IV.2.3.1.1-def-fr}
Sean $X$ un preesquema, $\mathcal F$ un $\mathcal O_X$-módulo
cuasicoherente. Se dice que un punto $x\in X$ está asociado a $\mathcal F$ si
el ideal maximal $\mathfrak m_x$ de $\mathcal O_x$ está asociado al
$\mathcal O_x$-módulo $\mathcal F_x$ \textup{(es decir, si $\mathfrak m_x$ es el
\emph{anulador} de un elemento de $\mathcal F_x$)}. Se denota por
$\operatorname{Ass}(\mathcal F)$ el conjunto de los $x\in X$ asociados a
$\mathcal F$. Se dice que una parte cerrada irreducible $Z$ de $X$ es un
ciclo primo asociado a $\mathcal F$ si su punto genérico está
asociado a $\mathcal F$. Cuando $\mathcal F=\mathcal O_X$, se dirá también que
los puntos \textup{(resp.} ciclos primos) asociados a $\mathcal F$ están
asociados al preesquema $X$.
\end{definition}

Se dice que un ciclo primo asociado a $\mathcal F$ (resp. $X$) es
\emph{inmerso} si está contenido en \emph{otro} ciclo primo asociado a
$\mathcal F$ (resp. $X$) (es decir, si no es \emph{maximal} en el conjunto de
estos ciclos).

Si $X$ es localmente noetheriano, las \emph{componentes irreducibles} de $X$
no son sino los ciclos primos \emph{maximales} (o \emph{no inmersos})
asociados a $X$.

Es claro que si $x\in\operatorname{Ass}(\mathcal F)$, se tiene
$\mathcal F_x\ne0$, es decir,

\begin{env}[3.1.1.1]
\label{IV.2.3.1.1.1-fr}
\[
\operatorname{Ass}(\mathcal F)\subset\operatorname{Supp}(\mathcal F).
\]
\end{env}

Si $x\in X$ está asociado a $\mathcal F$, está evidentemente asociado a
$\mathcal F|U$ para toda vecindad abierta $U$ de $x$ y, recíprocamente, si está
asociado a $\mathcal F|U$ para una de estas vecindades, está asociado a
$\mathcal F$.

Se observará, por último, que para un $\mathcal O_X$-módulo cuasicoherente
$\mathcal F$, la existencia de ciclos primos asociados inmersos es una cuestión
\emph{local}, pues si $y$ y $z$ son dos puntos de
$\operatorname{Ass}(\mathcal F)$ tales que $y\in\overline{\{z\}}$, toda vecindad
de $y$ contiene a $z$.

